\documentclass[11pt,a4paper]{article}
% Spolecny styl pro ucitelske pripravy MAT (Work/MAT)
% Verze 1.1, 2026-10-23
\NeedsTeXFormat{LaTeX2e}
\ProvidesPackage{priprava}[2026/10/23 v1.1 Sablona ucitelskych priprav MAT]

% Zakladni typografie a cestina
\RequirePackage[T1]{fontenc}
\RequirePackage[utf8]{inputenc}
\RequirePackage[czech]{babel}
\RequirePackage{lmodern}
\RequirePackage{microtype}
\RequirePackage{amsmath,amssymb}
\RequirePackage{xcolor}
\RequirePackage{enumitem}
\RequirePackage[most]{tcolorbox}
\RequirePackage{tabularx,booktabs,array}
\RequirePackage{fancyhdr}
\RequirePackage{needspace}
\RequirePackage{geometry}
\RequirePackage[hidelinks]{hyperref}

\geometry{a4paper,top=17mm,bottom=18mm,left=23mm,right=23mm}
\definecolor{notegray}{gray}{0.39}
\setlength{\parindent}{0pt}
\setlength{\parskip}{2.5pt plus 1pt minus 0.5pt}
\setlist[itemize]{topsep=2pt,itemsep=1pt,parsep=0pt,leftmargin=1.45em}
\setlist[enumerate]{topsep=2pt,itemsep=2pt,parsep=0pt,leftmargin=1.65em}
\setcounter{secnumdepth}{0}
\renewcommand{\arraystretch}{1.12}

% Zahlavi: vlevo rocnik | globalni tema, vpravo cislo: nazev hodiny
% Zapati zustava zcela prazdne.
\pagestyle{fancy}
\fancyhf{}
\renewcommand{\headrulewidth}{0.35pt}
\renewcommand{\footrulewidth}{0pt}
\setlength{\headheight}{14pt}
\newcommand{\nastavPripravu}[4]{%
  \fancyhead[L]{\small #1\enspace\textbar\enspace #2}%
  \fancyhead[R]{\small #3: #4}%
}

% Minimalisticke cernobile boxy. Kazda definice zvlast.
\tcbset{commonlesson/.style={enhanced,breakable=false,colback=white,colframe=black,
  colbacktitle=white,coltitle=black,boxrule=0.45pt,
  arc=0.8mm,outer arc=0.8mm,left=3mm,right=3mm,top=1.6mm,bottom=1.6mm,
  toptitle=1mm,bottomtitle=0.7mm,fonttitle=\bfseries,
  attach boxed title to top left={xshift=2.5mm,yshift=-1.15mm},
  boxed title style={colback=white,colframe=white,boxrule=0pt,left=1mm,right=1mm}}}
\newtcolorbox{zadaniBox}[1][]{commonlesson,title={Zadání},#1}
\newtcolorbox{prikladBox}[1][]{commonlesson,title={Příklad},#1}
\newtcolorbox{definiceBox}[1][]{commonlesson,title={Definice},#1}

% Spolecny box pro zadani a jeho reseni, bez mezery mezi castmi.
% Pouziti:
%   \begin{prikladReseniBox}
%     Zadani ulohy.
%     \tcblower
%     \begin{align*} ... \end{align*}
%   \end{prikladReseniBox}
% Spodni cast ma jemny sedy podklad a delici linku; lze v ni sazet
% vice radku matematiky. Obe casti maji jen jeden spolecny ramecek.
\tcbset{solutionlesson/.style={
  commonlesson,
  enhanced jigsaw,
  breakable,
  colbacklower=black!3,
  segmentation style={solid,draw=black!25,line width=0.35pt},
  fontlower=\small\color{notegray},
  before lower={\noindent\textbf{Řešení:}\enspace},
  leftlower=3mm,rightlower=3mm
}}
\newtcolorbox{prikladReseniBox}[1][]{solutionlesson,title={Příklad},#1}
\newtcolorbox{zadaniReseniBox}[1][]{solutionlesson,title={Zadání},#1}

% Pomucky pro autora pripravy: sedive poznamky; \reseni pro starsi pripravy
\newcommand{\poznamka}[1]{\par\begingroup\small\color{notegray}\itshape
  \textnormal{\textbf{Poznámka pro učitele.}}\enspace #1\par\endgroup}
\newcommand{\reseni}[1]{\par\begingroup\small\color{notegray}
  \textbf{Očekávané řešení:}\enspace #1\par\endgroup}
\newcommand{\blok}[1]{\needspace{5\baselineskip}\subsection*{#1}}
\newcommand{\vyrok}[1]{\uv{#1}}

\nastavPripravu{1. ročník}{Algebraické výrazy}{43}{Násobení jednočlenem}
\begin{document}
\poznamka{45 minut: 5 min aktivace distributivity, 10 min násobení s důkazem, 10 min společné příklady, 12 min samostatně, 5 min kontrola, 3 min zadání domácí práce. Součin dvou víceročlenných závorek systematicky řeší až hodina 44.}
\blok{Od číselné distributivity k mnohočlenům}
$3(4+5)=3\cdot4+3\cdot5$. Toto pravidlo platí i pro výrazy: každý člen v závorce vynásobíme činitelem před závorkou.
\begin{definiceBox}[title={Pravidlo: násobení jednočlenu mnohočlenem}]
Pro libovolné výrazy $a,u,v,w$ platí \textbf{distributivní zákon} $a(u+v+w)=au+av+aw$. Pro jednočleny zároveň násobíme koeficienty a sčítáme exponenty stejných proměnných: $(ax^m)(bx^n)=abx^{m+n}$.
\end{definiceBox}
\poznamka{Zdůrazni, že násobení \uv{přechází} ke každému členu; nepoužij pravidlo $a(u+v)=au+v$. Pracujeme pouze s nezápornými celočíselnými exponenty proměnných.}
\blok{Řešené příklady}
\begin{prikladReseniBox}[title={Příklad 1: kladný jednočlen}]
Roznásob $3x(2x^2-4x+5)$.
\tcblower
$3x\cdot2x^2+3x\cdot(-4x)+3x\cdot5=\boxed{6x^3-12x^2+15x}$.
\end{prikladReseniBox}
\begin{prikladReseniBox}[title={Příklad 2: záporný jednočlen}]
Uprav $-2a^2(3a^3-4a+7)$.
\tcblower
$-6a^5+8a^3-14a^2$. Znaménko minus musí být součástí násobení každého členu.
\end{prikladReseniBox}
\begin{prikladReseniBox}[title={Příklad 3: více proměnných}]
Uprav $-3xy(2x^2y-4xy^2+5)$.
\tcblower
$\boxed{-6x^3y^2+12x^2y^3-15xy}$. Při násobení $xy\cdot x^2y$ se exponenty obou proměnných sčítají zvlášť.
\end{prikladReseniBox}
\blok{Samostatné procvičování}
\begin{zadaniReseniBox}[title={Samostatná práce 1: od jednoduchého ke složitějšímu}]
\begin{enumerate}[label=\alph*)]
\item $4x(x^2-3x+2)$.
\item $-5a(2a^2-a-1)$.
\item $3x^2y(2xy-y^2+4)$.
\item $-\frac12 x(6x^2-4x+8)$.
\end{enumerate}
\tcblower
\begin{enumerate}[label=\alph*)]
\item $4x^3-12x^2+8x$.
\item $-10a^3+5a^2+5a$.
\item $6x^3y^2-3x^2y^3+12x^2y$.
\item $-3x^3+2x^2-4x$.
\end{enumerate}
\end{zadaniReseniBox}
\begin{zadaniReseniBox}[title={Samostatná práce 2: násobení a slučování}]
\begin{enumerate}[label=\alph*)]
\item $2x(x^2+3x-1)-x(x^2-x+2)$.
\item $-2a(3a-4)+3a(a+1)$.
\item $3x(2x-1)-(x^2+4x-5)$.
\end{enumerate}
\tcblower
\begin{enumerate}[label=\alph*)]
\item $x^3+7x^2-4x$.
\item $-3a^2+11a$.
\item $5x^2-7x+5$.
\end{enumerate}
\end{zadaniReseniBox}
\blok{Kde se to používá?}
\begin{prikladReseniBox}[title={Příklad 4: obsah obdélníku}]
Obdélník má délky stran $3x$ a $(2x+5)$ centimetrů, $x>0$. Zapiš obsah mnohočlenem a urči jej pro $x=2$.
\tcblower
$S=3x(2x+5)=6x^2+15x\,\mathrm{cm}^2$; po dosazení $S(2)=24+30=\boxed{54\,\mathrm{cm}^2}$.
\end{prikladReseniBox}
\begin{zadaniReseniBox}[title={Rozšiřující příklad 1: zpětné určení výrazu}]
Najdi mnohočlen $P(x)$ takový, aby $-3xP(x)=12x^4-6x^3+9x^2$ pro všechna $x$.
\tcblower
$P(x)=\boxed{-4x^3+2x^2-3x}$. Kontrola roznásobením. Závěr plyne z rovnosti koeficientů mnohočlenů.
\end{zadaniReseniBox}
\blok{Domácí procvičení}
\begin{zadaniReseniBox}[title={Domácí úkol 1}]
\begin{enumerate}[label=\alph*)]
\item $7x(2x^2-x+3)$.
\item $-2ab(3a^2-2ab+5b^2)$.
\item $\frac13 x^2(9x^3-6x+12)$.
\item $4a(a^2-3a+2)-2a(a^2+a-1)$.
\item Vypočti obsah obdélníku se stranami $5t$ a $t+4$ pro $t=3$, $t>0$.
\end{enumerate}
\tcblower
\begin{enumerate}[label=\alph*)]
\item $14x^3-7x^2+21x$.
\item $-6a^3b+4a^2b^2-10ab^3$.
\item $3x^5-2x^3+4x^2$.
\item $2a^3-14a^2+10a$.
\item $S=5t^2+20t$, $S(3)=45+60=105$ (plošných jednotek).
\end{enumerate}
\end{zadaniReseniBox}
\end{document}
